This study shows that modifying the training objective is an effective way
to control model behavior in sparse lesion segmentation. Rather than
increasing architectural complexity, the proposed approach redistributes
supervision during optimization, enabling clearer interpretation of its
effects. Voxel-wise objectives inherently favor high-volume regions, causing small
lesions to have limited influence. The proposed objective addresses this
through two terms: a component-adaptive term that balances contributions
across lesion structures, and a lesion-level term that enforces lesion
detection. Together, these shift learning from purely overlap-driven
optimization toward better sensitivity to sparse lesion signals. Results support this design. The method improves small lesion recall and reduces false negatives while maintaining competitive Dice and improving
boundary accuracy, indicating balanced performance across metrics. However, increased sensitivity leads to reduced lesion-wise precision, often due to small isolated false positives or fragmented predictions.
This reflects a trade-off between detection and delineation, further
affected by the lesion-level term, which enforces detection but not full
coverage. Lesion detectability also depends on factors such as contrast and boundary
clarity. While the proposed objective improves sensitivity, it does not
fully overcome limitations in input data. Limitations include evaluation on a single dataset, focus on one disease,
and use of a single architecture. In addition, the detection–precision
trade-off is not explicitly controlled. Future work should evaluate generalization across datasets and diseases,
validate across architectures, and improve control over false positives
and boundary consistency. Extending lesion-level supervision to include
coverage constraints may further improve segmentation quality. Overall, the results suggest that small lesion segmentation depends on how sparse signals are represented during learning. The proposed objective
improves sensitivity to small lesions while maintaining stable global
performance, highlighting the role of loss design in imbalanced settings.